% 第4章第7节：本章小结。

\section{本章小结}
\label{sec:graph-set-functions-summary}

本章从传感器覆盖出发，把对整组选择的评价写成有限底集上的集合函数。元素之间若互不影响，每个元素贡献一个固定的数，集合函数就是模函数；覆盖范围的重叠使边际增益随已选集合增大而递减，这就是次模性；部件之间的协同则产生相反方向的超模性。单调、模、次模各自约束不同的东西，覆盖量、割权和净收益说明它们可以以各种方式组合，判断时要回到边际增益或四个函数值的不等式逐项核对。子集累计与Möbius反演完成的是另一项任务：已知全部子集的累计值，按元素个数之差交替加减，恢复每个子集自身的贡献；它只依赖有限求和与包含关系，结果可以为负，计算开销随底集大小指数增长。

选择问题按目标方向与约束分别处理。在基数约束下最大化归一化、单调、次模的函数，逐轮加入当前增益最大的元素，所得值不低于最优值的$1-(1-1/k)^k$倍；撤去次模性，贪心会错过互补组合，撤去单调性，强制选满会降低目标。换成预算约束后，单位成本规则在传感器实例中找到了最优解，但在成本为$(1,2)$与$(100,100)$的例子中可以任意差，最佳单元素比较只修复个别失效。最小化由一元项和满足次模条件的成对项组成的函数时，目标可以表示为$s$--$t$网络的割容量，一次最小割给出精确解，并附带可以独立核对的证书。使用任何一种方法之前，都要同时写明目标方向、约束、单调性、次模性以及函数的表示结构。

第~\ref{sec:graph-pairwise-binary-cut}节的成功依赖两点：每个变量只有两种状态，成对项满足特定的不等式。许多问题两点都不满足，变量可以有更多取值，局部成本可以同时读取三个变量，或者直接禁止某些组合；但整体函数常常仍是许多局部项之和，每一项只读取少数变量。第~\ref{chap:graph-dynamic-programming}章从这里出发，研究怎样借助记录局部项读取关系的依赖图，把对全部配置的枚举化为沿图逐步传递的小规模表格计算。
